\section{Engine Architecture}
\label{sec:architecture}

This section describes the overall architecture of the implemented chess engine.
The focus lies on the structural design decisions and the interaction between the
core components. Rather than optimizing for maximum playing strength, the
architecture is deliberately kept modular and transparent in order to facilitate
experimentation and detailed analysis of search algorithms and time management
strategies.

\subsection{Overall Design}

The engine follows a layered architecture that separates game representation,
move generation, search, evaluation, and search control. Each component has a
clearly defined responsibility and communicates with other modules through
narrow and explicit interfaces.

\cref{fig:architecture} provides a high-level overview of the engine structure
and the interaction between the individual modules.

\begin{figure}[ht]
\centering
\begin{tikzpicture}[
  box/.style={draw, rounded corners, inner sep=6pt, align=center},
  arr/.style={-Latex, thick}
]
\node[box] (board) {Board State\\(Mailbox)};
\node[box, right=1cm of board] (movegen) {Move Generation};
\node[box, right=1cm of movegen] (search) {Search\\(ID + NM + Q)};
\node[box, below=1cm of search] (eval) {Evaluation\\(MG/EG + PSQT)};
\node[box, above=1cm of search] (tt) {Position Cache\\(TT)};
\node[box, below=1cm of movegen] (time) {Search Constraints\\(Time / Nodes / Depth)};
\node[box, right=1cm of search] (heuristics) {Heuristics - Move ord.\\(PV/TT+HH+KM+Material)};

\draw[arr] (board) -- (movegen);
\draw[arr] (movegen) -- (search);
\draw[arr] (search) -- (eval);
\draw[arr] (tt) -- (search);
\draw[arr] (time) -- (search);
\draw[arr] (heuristics) -- (search);
\end{tikzpicture}
\caption{High-level architecture of the implemented chess engine.}
\label{fig:architecture}
\end{figure}

At a conceptual level, the engine consists of the representation of the game
state, the generation and application of moves, the search procedure, the static
evaluation, and a unified constraint mechanism that controls search termination.
This separation ensures that individual components can be modified or replaced
without affecting the rest of the system, which is particularly important for
experimental analysis.

\subsection{Game State and Move Handling}

\paragraph{Board Representation}
The board is represented using a classical mailbox model consisting of 64
squares. Each square explicitly stores the occupying piece or indicates that the
square is empty. In addition to the piece placement, the board state contains
information about the side to move, castling rights, king squares and en-passant availability.

Modern state-of-the-art engines typically rely on bitboard-based representations
for efficient move generation and evaluation. However, mailbox representations
remain a common choice in educational and experimental engines due to their
simplicity and transparency~\cite{cpw_bitboards}. Given the focus of this work on
algorithmic behavior rather than raw performance, the mailbox model provides a
clear and easily analyzable foundation.

\paragraph{Move Representation}
Moves are represented as explicit objects containing all information required to
apply and undo them. This includes the source and destination squares, the moving
piece, captured pieces, promotion information, and special move flags for
castling, promotion and en-passant captures.

This design enables a fully reversible make--unmake mechanism, which is essential
for recursive search algorithms. By storing all relevant information directly in
the move object, the engine avoids recomputation when restoring previous board
states, thereby simplifying the search implementation and reducing potential
sources of errors.

\paragraph{Move Generation}
Move generation follows a direct and exhaustive approach. For each piece on the
board, all potential target squares are examined according to the rules of
chess. No precomputed attack tables or bitboard-based optimizations are used.

Illegal moves, such as those that leave the king in check, are filtered during
move execution rather than during generation. This design keeps the move
generator independent of check detection logic and makes its behavior easier to
reason about. Although this approach is less optimized than state-of-the-art
alternatives, it provides a robust and well-understood baseline for search
experimentation.

\subsection{Position Hashing and Transposition Tables}

\paragraph{Zobrist Hashing}
To uniquely identify board positions, the engine employs Zobrist hashing. Each
combination of square and piece type is associated with a random 64-bit value,
and the hash of a position is computed by XOR-combining the values corresponding
to the pieces currently on the board. Additional state information, such as the
side to move, castling rights, and en-passant availability, is also incorporated
into the hash.

Zobrist hashing is commonly used in game tree search because it allows board
positions to be mapped efficiently to fixed-size keys. While the technique
supports incremental updates during make--unmake operations, the implementation
used in this engine recomputes the hash from the board state directly. The technique itself was originally introduced by Zobrist~\cite{zobrist1970hashing} and has since become a standard component of modern chess engines.

\paragraph{Transposition Table Usage}
The computed hash values are used as keys in a transposition table, which stores
previously evaluated positions. Each entry contains the evaluation score, the
search depth at which the position was evaluated, a bound type (exact, lower
bound, or upper bound), and the best move found.

The transposition table serves two closely related purposes. First, it reduces
search redundancy by reusing results for positions that can be reached via
different move orders, which is common in chess. Without such a mechanism,
identical positions would be evaluated repeatedly, leading to a significant
increase in search effort. Second, transposition table entries improve move
ordering by prioritizing moves that proved effective in previous searches. A
simple aging mechanism ensures that outdated entries are gradually replaced,
preventing the table from being dominated by stale information. For practical
implementations and further discussion, see for example Warnock~\cite{warnock1988search}.

\subsection{Search Control and Integration}

\paragraph{Unified Search Constraints}
All termination conditions for a search invocation are encapsulated in a single
constraint structure. This includes depth limits, node limits, time budgets, and
special modes such as infinite search or pondering.

Encapsulating all stopping criteria in a unified constraint object avoids
scattered stop checks throughout the search code and ensures consistent behavior
across different search modes. This design is particularly important for
experimental analysis, as it guarantees that differences in search behavior are
caused by the chosen constraints rather than by implementation artifacts.

\paragraph{Engine Integration}
The engine core is coordinated by a game management component that maintains the
current position, applies and reverts moves, and invokes the search procedure.
This component acts as a bridge between the external UCI interface and the
internal search and evaluation modules, ensuring a clean separation between
protocol handling and algorithmic logic.

\subsection{Design Rationale}

The architectural decisions made in this engine deliberately prioritize clarity,
correctness, and experimental flexibility over maximum playing strength. Many
state-of-the-art engines employ highly optimized data structures and aggressive
micro-optimizations, which can obscure the underlying algorithmic behavior.

By choosing a modular and explicit design, this engine allows the effects of
search algorithms, heuristics, and time management strategies to be studied in
isolation. This aligns with the primary goal of this work, which is to analyze
the behavior of iterative deepening search under time constraints rather than to
compete directly with top-performing chess engines.